\section{Transport model and comparison of design problems}
\label{sec:model}

The wall pattern is \emph{quenched}: it is sampled once and remains fixed while
particles move. Expectations over patterns below describe variation among
channels, not higher moments of the displacement of one particle. The flow
and equilibrium affinity are fixed while the exchange rate and tangential
surface diffusivity may vary along the wall.

\subsection{Physical model and units}
Let $\Omega\subset\R^2$ be a bounded connected cross-section with $C^2$
boundary $\Gamma$. Assume that $\Gamma$ is connected; its arclength coordinate
is periodic with period $P$, and $A=|\Omega|$. The bulk transverse diffusivity
$D_b>0$, affinity $K>0$, and longitudinal velocity $u\in L^2(\Omega)$ are
fixed. Adsorbed particles have zero longitudinal advective velocity. For a
realization $\omega$, the desorption rate is $k=k_\omega\ge0$, the adsorption
coefficient is $Kk$, and $D(s)\ge0$ is tangential surface diffusivity. Our
uniform estimates assume
\begin{equation}\label{eq:rate-assumptions}
 0\le k\le k_{\max},\qquad \int_\Gamma k\dd s\ge\kappa>0.
\end{equation}

When axial molecular diffusion is included, assume
$0\le D_b^x<\infty$ and $D_s^x\in L^1(\Gamma)$ with $D_s^x\ge0$.
For smooth coefficients, bulk concentration $p(x,y,t)$ and wall
concentration $\rho(x,s,t)$ satisfy
\begin{align}
 \partial_t p+u(y)\partial_xp
 &=D_b\Delta_y p+D_b^x\partial_x^2p,\notag\\
 -D_b\partial_n p&=Kkp_\Gamma-k\rho,\label{eq:physical-pde}\\
 \partial_t\rho&=\partial_s(D\partial_s\rho)
       +D_s^x(s)\partial_x^2\rho+Kkp_\Gamma-k\rho.\notag
\end{align}
Here $p_\Gamma$ is the bulk trace and $n$ is the outward normal. The boundary
loss in the bulk equals the wall gain. The divergence form of surface
motion is part of the model; replacing it by $D\partial_s^2$ would generally
change the stationary measure. The rough-coefficient interpretation is
given below.

Set
\begin{equation}\label{eq:equilibrium}
 Z=A+KP,\qquad
 \pi(\dd y,\mathrm{bulk})=\frac{\dd y}{Z},\quad
 \pi(\dd s,\mathrm{wall})=\frac{K\dd s}{Z},\qquad
 V=\frac{1}{Z}\int_\Omega u\dd y.
\end{equation}
Indeed, constant bulk concentration and wall concentration $K$ times that
constant make every transverse flux vanish. The transverse dynamics are
reversible with respect to $\pi$, as the symmetric form below confirms.
Thus the stationary mean speed $V$ is independent of $k$ and $D$.
We assume $V\ne0$ when comparing divergent optimal values and write
\begin{equation}\label{eq:chi}
 \chi=\frac{KV^2}{Z}>0.
\end{equation}

Before taking a small-budget limit, all quantities must be expressed in
fixed units. Dimensionally, $[k]=T^{-1}$, $[D]=L^2T^{-1}$,
$[K]=L$, and the budget $M=\int_\Gamma D\dd s$ has units $L^3T^{-1}$.
Choose $\ell_{\rm ref}=P_{\rm phys}/(2\pi)$ and a fixed reference rate
$k_{\rm ref}>0$. Put
\begin{equation}\label{eq:units}
 s'=s_{\rm phys}/\ell_{\rm ref},\quad t'=k_{\rm ref}t_{\rm phys},\quad
 k'=k_{\rm phys}/k_{\rm ref},\quad
 D'=\frac{D_{\rm phys}}{k_{\rm ref}\ell_{\rm ref}^2},\quad
 M'=\frac{M_{\rm phys}}{k_{\rm ref}\ell_{\rm ref}^3}.
\end{equation}
Bulk lengths and $K$ are divided by $\ell_{\rm ref}$, velocities by
$k_{\rm ref}\ell_{\rm ref}$, and diffusivities by
$k_{\rm ref}\ell_{\rm ref}^2$. We henceforth suppress primes and use fixed
dimensionless units. We retain $P$ in the comparison lemmas to show their
geometric dependence; the canonical ensemble has $P=2\pi$. The surface
response defined below has physical units $LT$ and converts as
\begin{equation}\label{eq:J-dimensional}
 J_{\rm phys}=\frac{\ell_{\rm ref}}{k_{\rm ref}}J',\qquad
 D_{\flow,\rm phys}=k_{\rm ref}\ell_{\rm ref}^2D'_{\flow}.
\end{equation}
In particular, a dimensionless scalar asymptotic is multiplied by
$\chi_{\rm phys}\ell_{\rm ref}/k_{\rm ref}$ when used for physical
flow-induced diffusivity. Every subsequent power or logarithm of $M$
refers to the dimensionless budget.

\subsection{Scalar response and admissible transverse dynamics}
\label{sec:forms}
For \emph{every} nonnegative $D\in L^1(\Gamma)$ define the smooth-test
functional
\begin{align}
 \mathfrak a_{k,D}[v]
   &=\int_\Gamma D|v'|^2\dd s+\int_\Gamma kv^2\dd s,
       &&v\in C^\infty(\Gamma),\label{eq:scalar-form}\\
 J_k(D)&=\sup_{v\in C^\infty(\Gamma)}
       \left\{2\int_\Gamma v\dd s-\mathfrak a_{k,D}[v]\right\}
       \in(0,\infty].\label{eq:J-definition}
\end{align}
Optimization over the scalar multiple of each test gives the equivalent
quotient formula
\begin{equation}\label{eq:J-quotient}
 J_k(D)=\sup_v\frac{(\int_\Gamma v\dd s)^2}{\mathfrak a_{k,D}[v]}.
\end{equation}
A zero denominator with nonzero numerator gives $+\infty$; the quotient
$0/0$ is assigned zero. These definitions impose no existence of a
minimizer, inverse operator, or diffusion associated with an arbitrary
measurable weight.

That distinction matters: a weighted derivative preform need not be
closable for every nonnegative integrable weight; the one-dimensional
closability criterion is Hamza's condition
\citep[Theorem~3.1.6]{FukushimaOshimaTakeda2011}. We call a coefficient
\emph{physically admissible} if
\begin{equation}\label{eq:surface-preform}
 e_D[v]=\int_\Gamma D|v'|^2\dd s,\qquad v\in C^\infty(\Gamma),
\end{equation}
is closable in $L^2(\Gamma)$. Explicitly, closability means that for
every smooth sequence $v_n$ with $v_n\to0$ in $L^2$ and
$e_D[v_n-v_m]\to0$ as $n,m\to\infty$, one has $e_D[v_n]\to0$.
This condition ensures that the form-norm completion embeds
unambiguously in $L^2$: two energy-Cauchy approximations to the same
$L^2$ function define the same energy and hence a unique minimal closed
form. For this class we use its minimal closure,
with domain $\mathcal V_D$, rather than choosing unspecified interface
conditions at zero-mobility sets. The closure norm is
$(\|v\|_2^2+e_D[v])^{1/2}$. The closed form
$\mathfrak a_{k,D}=e_D+\int kv^2$ has the same domain because $k$ is
bounded. No assertion of closability is needed for the larger scalar
optimization class in \eqref{eq:J-definition}.

\begin{lemma}[Positive-background realization]\label{lem:positive-form}
If $D\in L^1(\Gamma)$ and $D\ge m>0$, then \eqref{eq:surface-preform}
is closable, $\mathcal V_D\subset H^1(\Gamma)$, and, under
\eqref{eq:rate-assumptions},
\begin{equation}\label{eq:anchored}
 \|v\|_2^2\le C\left(\|v'\|_2^2+\int_\Gamma kv^2\dd s\right),
 \qquad
 \mathfrak a_{k,D}[v]\ge c\min(m,1)\|v\|_2^2.
\end{equation}
Consequently $h=H_{k,D}^{-1}1\in\mathcal V_D$ exists uniquely, where
$H_{k,D}$ is the self-adjoint operator of the closed form, and
\begin{equation}\label{eq:h-identities}
 h\ge0,\qquad J_k(D)=\int_\Gamma h\dd s
       =\mathfrak a_{k,D}[h],\qquad \int_\Gamma kh\dd s=P.
\end{equation}
All constants depend only on $P,k_{\max},\kappa$.
\end{lemma}
\begin{proof}
Suppose $v_n\to0$ in $L^2$ and $\sqrt D\,v_n'$ is Cauchy in $L^2$.
The floor implies that $v_n'$ is Cauchy in $L^2$, with limit zero by
distributional differentiation. A subsequence converges pointwise almost
everywhere both before and after multiplication by $\sqrt D$; since
$D$ is finite almost everywhere, the latter limit is also zero. This is
closability, and the same argument identifies the derivative of any
member of the closure with its weak $H^1$ derivative.

Write $\bar v=P^{-1}\int v$. Periodic Poincar\'e and the rate anchor give
\[
 \kappa|\bar v|^2
 \le 2\int_\Gamma kv^2+2k_{\max}\|v-\bar v\|_2^2
 \le2\int_\Gamma kv^2+C\|v'\|_2^2.
\]
Adding the mean-zero part proves \eqref{eq:anchored}, first for smooth
functions and then by closure. The closed-form representation theorem
and this bound give the inverse and its weak equation
$\mathfrak a_{k,D}(h,v)=\int v$.
The closure is stable under normal contractions: for a smooth
approximation to a Lipschitz contraction $T$ with $T(0)=0$ and
$|T'|\le1$, the chain rule gives $e_D[T(v)]\le e_D[v]$ on the smooth
core; approximation and weak lower semicontinuity extend the inequality
to the closure. In particular $h^-\in\mathcal V_D$, and testing with
$h^-$ gives $-\mathfrak a_{k,D}[h^-]=\int h^-\ge0$, so $h^-=0$.
Testing with $h$ and with the constant $1$ proves the identities.
The smooth core is dense in the form norm, so its supremum
\eqref{eq:J-definition} equals the inverse response just obtained.
\end{proof}

For a physically admissible $D$, the unnormalized transverse energy on
$H^1(\Omega)\oplus\mathcal V_D$ is
\begin{equation}\label{eq:coupled-energy}
 \mathcal E_D[f,v]=D_b\int_\Omega|\nabla f|^2\dd y
       +K e_D[v]+K\int_\Gamma k(v-f_\Gamma)^2\dd s.
\end{equation}
The trace map $H^1(\Omega)\to H^{1/2}(\Gamma)\subset L^2(\Gamma)$
is bounded. Thus the exchange term is continuous in the product form
norm; equivalently, a sequence Cauchy for $\mathcal E_D+\|\cdot\|_{L^2(\pi)}^2$
is Cauchy in the bulk $H^1$ and surface form norms and its exchange term
converges. This proves closedness. The same contraction argument,
including $|T(v)-T(f_\Gamma)|\le|v-f_\Gamma|$, proves the Markov property.
Smooth pairs are a form core and are uniformly dense in the continuous
functions on the disjoint union $\overline\Omega\sqcup\Gamma$.
Consequently this is a regular Dirichlet form. Its normalized version
$\mathcal E_D/Z$ defines a conservative symmetric Markov process with
stationary probability $\pi$; constants have zero energy. We use the
standard regular-form/process correspondence in this step
\citep[Theorem~7.2.1]{FukushimaOshimaTakeda2011}. For smooth coefficients,
integration by parts gives bulk generator $D_b\Delta f$, wall generator
$\partial_s(Dv')+k(f_\Gamma-v)$, and boundary condition
$D_b\partial_n f=Kk(v-f_\Gamma)$, recovering the transverse part of
\eqref{eq:physical-pde}.

If $D\ge m>0$, this process has a spectral gap and its stationary
law $\pi$ is ergodic. To check the needed coercivity, use the gauge
$\int_\Omega f=0$, obtainable by subtracting a common constant from
$(f,v)$. Bulk Poincar\'e and trace estimates bound $\|f\|_{H^1}$ and
$\|f_\Gamma\|_2$ by $C\|\nabla f\|_2$. Moreover,
\[
 \int_\Gamma kv^2
 \le2\int_\Gamma k(v-f_\Gamma)^2+2k_{\max}\|f_\Gamma\|_2^2.
\]
Together with \eqref{eq:anchored}, this bounds the squared product $L^2$
norm in this gauge by $C_m\mathcal E_D[f,v]$. Subtracting the
$\pi$-mean can only decrease the $L^2(\pi)$ norm. Hence the energy
controls the norm on the mean-zero subspace, which proves the gap and
that the kernel consists exactly of constants. No such gap is asserted
for all degenerate designs.

\subsection{Stationary dispersion and the Schur identity}
\label{sec:schur}
Let $Y_t$ denote the stationary transverse process. Its axial advective
velocity is $u(y)$ in the bulk and zero on the wall. The centered
velocity $g$ therefore equals $u-V$ and $-V$, respectively, and belongs
to $L^2(\pi)$ with mean zero. Let $\mathscr A_D$ be the nonnegative
self-adjoint operator associated with $\mathcal E_D/Z$. For the
advective displacement $X_t^{\rm adv}=\int_0^t(g(Y_r)+V)\dd r$,
stationarity and the semigroup identity give
\begin{equation}\label{eq:green-kubo-finite}
 \frac{\Var X_t^{\rm adv}}{2t}
 =\int_0^t(1-r/t)\langle g,e^{-r\mathscr A_D}g\rangle_\pi\dd r.
\end{equation}
The spectral measure $\nu_g$ of $\mathscr A_D$ is nonnegative, so the
right-hand side increases to
\begin{equation}\label{eq:dispersion-spectral}
 D_\flow=\int_{[0,\infty)}\lambda^{-1}\nu_g(\dd\lambda)
 \in[0,\infty],
\end{equation}
where mass at zero gives $+\infty$. Indeed the scalar integrand before
the limit is $\int_0^t(1-r/t)e^{-\lambda r}\dd r$, increasing to
$1/\lambda$ for $\lambda>0$ and equal to $t/2$ at zero; monotone
convergence applies. This proves existence of the stationary extended
coefficient directly from particle displacement. It is finite under a
positive mobility floor by the spectral gap. Finiteness implies
$\Var X_t^{\rm adv}=2D_\flow t+o(t)$; no central limit theorem is needed
for this assertion.

Independent axial Brownian motion with local diffusivities $D_b^x$ and
$D_s^x(s)$ adds
\begin{equation}\label{eq:molecular}
 D_{\rm mol}=\frac{A D_b^x+K\int_\Gamma D_s^x(s)\dd s}{Z}.
\end{equation}
Here the axial Brownian driver is independent of the transverse process.
The stated diffusivity assumptions and stationarity give expected
integrated diffusivity $tD_{\rm mol}<\infty$ on every finite interval.
Thus the axial stochastic integral is square integrable, with variance
$2tD_{\rm mol}$. Moreover,
conditional on the transverse path, the axial stochastic integral has mean zero,
so there is no covariance with the advective displacement. If surface
diffusion is isotropic, $D_s^x=D$ and its contribution is $KM/Z$.
The optimization theorems below concern $D_\flow$.

Completing squares in the spectral representation, or first applying
$(\mathscr A_D+\epsilon)^{-1}$ and letting $\epsilon\downarrow0$,
gives
\begin{equation}\label{eq:full-variational}
 ZD_\flow=\sup_{f,v}\{2\mathcal F(f,v)-\mathcal E_D[f,v]\},\qquad
 \mathcal F(f,v)=\int_\Omega(u-V)f\dd y-KV\int_\Gamma v\dd s.
\end{equation}
The supremum is over the closed form domain, equivalently over smooth
pairs. For completeness, the general identity used here is
$\sup_w(2\langle g,w\rangle-\|\mathscr A_D^{1/2}w\|_\pi^2)
=\int\lambda^{-1}\dd\nu_g$. If the kernel projection
$g_0=\mathbf1_{\{0\}}(\mathscr A_D)g$ is nonzero, the trials
$w_n=ng_0$ have zero energy and objective
$2n\|g_0\|_\pi^2\to\infty$, agreeing with the spectral integral.
Otherwise, on the positive spectral subspace the upper inequality
follows from $2ab-\lambda b^2\le a^2/\lambda$. The form-domain trials
$w_n=\mathbf1_{[1/n,n]}(\mathscr A_D)\mathscr A_D^{-1}g$,
with the inverse interpreted only on the indicated spectral interval,
give objectives $\int_{[1/n,n]}\lambda^{-1}\dd\nu_g$.
Monotone convergence proves the reverse inequality, including a
divergent positive-spectrum integral. Thus \eqref{eq:full-variational}
is a consequence, not a
definition, of the physical coefficient.

\begin{proposition}[Surface lower bound and finite-response identity]
\label{prop:schur}
Every physically admissible design obeys
\begin{equation}\label{eq:bulk-lower}
 D_\flow(D,k)\ge\chi J_k(D),
\end{equation}
with extended values allowed. If $D\ge m>0$, let $h=H_{k,D}^{-1}1$
and $w=kh$. Then
\begin{equation}\label{eq:schur-identity}
 D_\flow(D,k)=\chi J_k(D)+\mathcal R(D,k),\qquad \mathcal R\ge0,
\end{equation}
where
\begin{align}
 Z\mathcal R&=\sup_{\int_\Omega f=0}
   \left\{2\mathcal L_D(f)-D_b\int_\Omega|\nabla f|^2-K\mathcal S_D(f_\Gamma)\right\},
       \label{eq:bulk-remainder}\\
 \mathcal L_D(f)&=\int_\Omega(u-V)f-KV\int_\Gamma wf_\Gamma,\notag\\
 \mathcal S_D(b)&=\inf_{v\in\mathcal V_D}
       \left\{e_D[v]+\int_\Gamma k(v-b)^2\right\}\ge0.\notag
\end{align}
The same identity holds for every physically admissible $D$ for which
$J_k(D)<\infty$, using the energy-space interpretation of $h$ and $w$
given in the proof. In that case $\mathcal R$ is finite, although no
uniform bound is claimed without a floor.
\end{proposition}
\begin{proof}
In \eqref{eq:full-variational}, take $f=0$, $v=-V\phi$, and then the
supremum over smooth $\phi$. This proves \eqref{eq:bulk-lower} without
assuming a finite inverse. For the identity, fix $f\in H^1(\Omega)$
and set $b=f_\Gamma$. The surface terms to be maximized are
\[
 K\{2\langle kb-V,v\rangle-\mathfrak a_{k,D}[v]-\int kb^2\}.
\]
Since $kb\in L^2$ and $H_{k,D}$ is coercive, their maximum is
$K\{\langle kb-V,H_{k,D}^{-1}(kb-V)\rangle-\int kb^2\}$.
Expand this expression. Its constant term is $KV^2J_k(D)$, its term
linear in $b$ is $-2KV\int wb$, and its quadratic term is
$-K\mathcal S_D(b)$, because completing the square also yields
$\mathcal S_D(b)=\int kb^2-\langle kb,H_{k,D}^{-1}kb\rangle$.
All quantities are finite. Finally
$\int_\Omega(u-V)=KPV$ and $\int w=P$, so
$\mathcal L_D$ annihilates constants, and $\mathcal S_D(b+a)=\mathcal S_D(b)$
by shifting its trial $v$ by $a$. The bulk mean-zero gauge is therefore
legitimate. Choosing $f=0$ proves $\mathcal R\ge0$.

For clarity, the extension to finite response without a floor does not
presume an $L^2$ corrector. Complete smooth functions, modulo zero-energy
functions, in the norm $\mathfrak a_{k,D}^{1/2}$ and call the resulting
Hilbert space $\mathcal W$. By \eqref{eq:J-quotient}, the source
$\ell(v)=\int v$ extends continuously to $\mathcal W$ exactly when
$J_k(D)<\infty$. Its Riesz representative $h\in\mathcal W$ satisfies
$\|h\|_{\mathfrak a}^2=J_k(D)$. The map $v\mapsto\sqrt k\,v$
is continuous from the smooth energy space to $L^2(\Gamma)$ and thus
extends to $\mathcal W$; write $w=\sqrt k(\sqrt k\,h)\in L^2$.
Here $\sqrt k\,h$ denotes that extended image, not an assumed pointwise
$L^2$ function $h$. Testing against the constant gives $\int w=P$.
For $b\in L^2(\Gamma)$, the functional
$\ell_b(v)=\int kbv$ is also continuous in this energy norm, since
$|\ell_b(v)|\le\|\sqrt k\,b\|_2\mathfrak a_{k,D}[v]^{1/2}$.
Let $z_b$ be its Riesz representative. The surface supremum is now
$K\{\|z_b-Vh\|_{\mathfrak a}^2-\int kb^2\}$, with cross term
$-2KV\int wb$. Moreover
$\mathcal S_D(b)=\int kb^2-\|z_b\|_{\mathfrak a}^2\ge0$;
the infimum over $\mathcal V_D$ equals the one over smooth tests and
then over $\mathcal W$, by their defining density. This proves the
same completion without subtracting infinite energies. Finally
$w\in L^2$, bulk Poincar\'e, and the trace estimate bound
$|\mathcal L_D(f)|\le C_{D,k}\|\nabla f\|_2$ in the mean-zero
gauge. The supremum in \eqref{eq:bulk-remainder} is therefore finite.
\end{proof}

\subsection{A logarithmic bound uniform in the mobility field}
\begin{theorem}[Finite-bulk remainder]\label{thm:log-remainder}
Under \eqref{eq:rate-assumptions}, if $D\in L^1(\Gamma)$ and
$D\ge m>0$, then
\begin{equation}\label{eq:log-remainder}
 0\le\mathcal R(D,k)\le C\{1+\log_+(1/m)\}.
\end{equation}
The constant depends only on the fixed geometry, $D_b>0$, $K$,
$\|u\|_{L^2(\Omega)}$, $k_{\max}$, and $\kappa$.
In particular it is independent of derivatives and upper bounds of $D$.
Here $\log_+x=\max(0,\log x)$ in the fixed units of \eqref{eq:units}.
\end{theorem}
\begin{proof}
Replace $m$ by $\min(m,1)$, so it suffices to take $0<m\le1$.
By \eqref{eq:anchored} and Cauchy--Schwarz,
$J=\langle1,H_{k,D}^{-1}1\rangle\le C/m$.
The energy identity gives $\|h'\|_2^2\le J/m\le C/m^2$.
Since $h\ge0$, its mean is $J/P$, and the fundamental theorem of
calculus for $H^1$ functions on the circle gives
\begin{equation}\label{eq:w-infty}
 \|h\|_\infty\le J/P+\sqrt P\|h'\|_2\le C/m,
 \qquad 0\le w=kh\le C/m,\quad \int_\Gamma w=P.
\end{equation}

We spell out the estimate that prevents this crude supremum bound from
producing a power-divergent bulk error. Normalize Fourier coefficients
by $\widehat w_n=P^{-1}\int_0^Pw(s)e^{-2\pi i ns/P}\dd s$ and use
\[
 \|z\|_{H^{-1/2}(\Gamma)}^2
   =P\sum_{n\in\mathbb Z}(1+(2\pi n/P)^2)^{-1/2}|\widehat z_n|^2.
\]
Positivity and the mass identity imply $\widehat w_0=1$,
$|\widehat w_n|\le1$, and Parseval gives
$\sum_n|\widehat w_n|^2=P^{-1}\|w\|_2^2\le\|w\|_\infty$.
For an integer $N\ge2$, separate the nonzero modes at $N$:
\begin{align}
 \|w-1\|_{H^{-1/2}}^2
 &\le C\sum_{n=1}^N\frac1n+\frac C N\sum_{|n|>N}|\widehat w_n|^2\notag\\
 &\le C\left(1+\log N+\frac{\|w\|_\infty}{N}\right).
\end{align}
Taking $N=\lceil2+\|w\|_\infty\rceil$ and using
\eqref{eq:w-infty} yields
\begin{equation}\label{eq:negative-half}
 \|w-1\|_{H^{-1/2}}^2\le C\{1+\log(1/m)\}.
\end{equation}

In the gauge $\int_\Omega f=0$, bulk Poincar\'e, the trace theorem,
and $H^{-1/2}$--$H^{1/2}$ duality imply
\[
 |\mathcal L_D(f)|
 \le C\{1+\|w-1\|_{H^{-1/2}}\}\|\nabla f\|_2.
\]
Indeed separate $w=1+(w-1)$ in \eqref{eq:bulk-remainder}; both the
bulk load and the integral of the trace are bounded by the bulk gradient,
and the remaining pairing is bounded by the two trace-dual norms.
Dropping $K\mathcal S_D\ge0$ and maximizing $2ax-D_bx^2$ over
$x\ge0$ gives $Z\mathcal R\le a^2/D_b$.
Equation \eqref{eq:negative-half} proves the claim.
\end{proof}

The connected one-dimensional wall is used both in the anchored
Poincar\'e estimate and in the logarithmic Fourier sum. A disconnected
wall would require additional componentwise anchors. This theorem is not
uniform as $D_b$ tends to zero, and it does not bound the remainder of
arbitrary designs without a positive floor.

\subsection{Information policies and transfer of optimal moments}
\label{sec:information-transfer}
Let $(\omega,Y)$ have a specified joint probability law on standard
Borel spaces. The observation $Y$ is the information available before
the mobility is chosen. The map $\omega\mapsto k_\omega$ is Borel as a
map into $L^1(\Gamma)$ and satisfies \eqref{eq:rate-assumptions} almost
surely. A scalar policy is a Borel map $y\mapsto D_y\in L^1(\Gamma)$
with $D_y\ge0$ and $\int D_y=M$ for every $y$.
A physical policy additionally uses physically admissible coefficients.
We impose no minimizer existence and no interchange of expectation and
pointwise infimum. Define
\begin{align}
 F_q(M;Y)&=\inf_{\text{scalar policies}}
                  \E[J_{k_\omega}(D_Y)^q],\label{eq:F-definition}\\
 G_q(M;Y)&=\inf_{\text{physical policies}}
                  \E[D_\flow(D_Y,k_\omega)^q],\qquad q>0.
                  \label{eq:G-definition}
\end{align}
Every budget is per observation, rather than an average over observations.
The joint observation law may depend on $M$.

These objectives are measurable. To see this without a measurable
selection theorem, choose a countable family of real trigonometric
polynomials with rational coefficients, dense in $C^1(\Gamma)$.
For each smooth $v$ the scalar expression in \eqref{eq:J-definition}
is continuous in $(k,D)$ in the $L^1\times L^1$ topology; the countable
family has the same supremum because all its integrands converge under
$C^1$ approximation. Thus $J$ is a jointly Borel extended function.
For the physical response, choose similarly a countable $C^1$-dense
family in the smooth bulk--wall core. Each expression in
\eqref{eq:full-variational} is continuous in $(k,D)$; density of the
core in every admissible form domain gives the same countable supremum.
This defines a Borel function on the entire scalar coefficient space
which agrees with the physical response whenever the form is closable.
In particular, no assertion that the set of closable weights is itself
Borel is required. The uniform policy is physical, and its positive-floor
coercivity bound, uniform over $k$, shows that both infima are finite.

\begin{theorem}[Same-budget, same-information comparison]
\label{thm:optimal-transfer}
Let $q>0$, $0<M<1/2$, and $0<\theta<1$. With
$L_\theta(M)=1+\log_+(P/(\theta M))$,
\begin{equation}\label{eq:transfer-lower}
 \chi^qF_q(M;Y)\le G_q(M;Y).
\end{equation}
There is a constant independent of the information law such that
\begin{align}
 G_q(M;Y)^{1/q}
 &\le\frac{\chi}{1-\theta}F_q(M;Y)^{1/q}+C L_\theta(M),
                  &&q\ge1,\label{eq:transfer-q-large}\\
 G_q(M;Y)
 &\le\left(\frac{\chi}{1-\theta}\right)^qF_q(M;Y)
                         +C^q L_\theta(M)^q,
                  &&0<q<1.\label{eq:transfer-q-small}
\end{align}
Consequently, along any sequence of budgets and observation laws for which
\begin{equation}\label{eq:transfer-growth}
 \frac{F_q(M;Y)}{[1+\log(1/M)]^q}\longrightarrow\infty,
\end{equation}
one has $G_q(M;Y)\sim\chi^qF_q(M;Y)$ as $M\downarrow0$.
\end{theorem}
\begin{proof}
The physical policy class is contained in the scalar class, so
\eqref{eq:bulk-lower} proves \eqref{eq:transfer-lower}. For any scalar
policy, even one using nonclosable weights, form the same-information
mixture
\begin{equation}\label{eq:mixture}
 D_y^\theta=(1-\theta)D_y+\theta M/P.
\end{equation}
This is Borel, has exactly budget $M$ for every $y$, and is physical by
\cref{lem:positive-form}. On smooth tests,
\[
 \mathfrak a_{k,D_y^\theta}
 =(1-\theta)\mathfrak a_{k,D_y}
       +\theta\mathfrak a_{k,M/P}
 \ge(1-\theta)\mathfrak a_{k,D_y}.
\]
The quotient formula implies
$J_k(D_y^\theta)\le J_k(D_y)/(1-\theta)$, including infinite values.
By \cref{thm:log-remainder}, realization by realization,
\begin{equation}\label{eq:mixture-physical}
 D_\flow(D_y^\theta,k)
 \le\frac{\chi}{1-\theta}J_k(D_y)+C L_\theta(M).
\end{equation}
For $q\ge1$, apply Minkowski's inequality to this bound. For $0<q<1$,
use $(a+b)^q\le a^q+b^q$ for nonnegative $a,b$. Take a sequence of
scalar policies whose expected responses approach $F_q$ and let its
approximation error tend to zero. This proves
\eqref{eq:transfer-q-large}--\eqref{eq:transfer-q-small} without an
optimal policy or any expectation--infimum interchange.
Finally take $\theta=M$. Then $L_\theta(M)=O(1+\log(1/M))$ and
$(1-\theta)^{-1}\to1$; divide the inequalities by the scalar optimum
and use \eqref{eq:transfer-growth}.
\end{proof}

\begin{corollary}[Uniformity over observations in the cosine ensemble]
\label{cor:all-information}
Let $\Gamma=\R/(2\pi\mathbb Z)$, $k_c(s)=(c+\cos s)^2$, and
$c$ be uniform on $[-2,2]$. For every observation law with this marginal,
every fixed $q>0$, and all sufficiently small $M$,
\begin{equation}\label{eq:all-info-lower}
 F_q(M;Y)\ge c_qM^{-q/5},
\end{equation}
and, uniformly over those observation laws,
\begin{equation}\label{eq:all-info-ratio}
 1\le\frac{G_q(M;Y)}{\chi^qF_q(M;Y)}\le
 \begin{cases}
 1+C_qM^{q/5}[\log(1/M)]^q,&0<q<1,\\
 1+C_qM^{1/5}\log(1/M),&q\ge1.
 \end{cases}
\end{equation}
\end{corollary}
\begin{proof}
For $|c|\le1/2$, select either root $r$ of $c+\cos r=0$ and fix a
nonnegative, nonzero smooth bump $\psi$ supported in $(-1,1)$.
For $\ell=M^{1/5}$ small, the periodic test
$v(s)=\psi((s-r)/\ell)$ has
$\int v=a_\psi\ell>0$. Since cosine is Lipschitz with constant one,
$k_c(s)\le|s-r|^2$ on its support. Every budget-$M$ field therefore
satisfies
\[
 \int D|v'|^2\le\|\psi'\|_\infty^2M/\ell^2,
 \qquad \int k_cv^2\le C_\psi\ell^3.
\]
The quotient \eqref{eq:J-quotient} gives
$J_c(D)\ge c\ell^2/(M/\ell^2+\ell^3)=cM^{-1/5}$.
This bound is pointwise in both $c$ and $D$; it applies to any policy,
regardless of what it observes. The event $|c|\le1/2$ has probability
$1/4$, proving \eqref{eq:all-info-lower}. In this ensemble $k_c\le9$
and $\int k_c=2\pi(c^2+1/2)\ge\pi$, so all comparison constants
are uniform in $c$. Substitute $\theta=M$ and
\eqref{eq:all-info-lower} into
\eqref{eq:transfer-q-large}--\eqref{eq:transfer-q-small}.
For $q\ge1$, raise the resulting ratio of $q$th roots to the fixed
power $q$; its excess is $O(M^{1/5}\log(1/M))$.
For $q<1$ the excess is $O(M^{q/5}\log(1/M)^q)$.
The mixture error $O(M)$ is smaller in both cases.
\end{proof}

The comparison concerns optimal \emph{values} at identical budgets and
information. It does not assert that an unregularized scalar optimizer
exists or has a small bulk remainder. Its upper bound instead constructs
physical policies with a vanishing fraction of the budget in a uniform
background. The same conclusions hold after adding any uniformly bounded
nonnegative molecular contribution, including \eqref{eq:molecular} with
fixed bulk molecular diffusion and isotropic wall diffusion. Additional
fabrication constraints that exclude \eqref{eq:mixture}, a vanishing
bulk diffusivity, a zero mean velocity, or a nonuniformly anchored rate
family require separate analysis.
